\subsection{Capped contained dyadic squares}

% Auxiliary finite geometry from scanner:templates, not the complete entropy lemma.
\begin{definition}[Lattice dyadic squares]
    \label{def:area_law_lattice_dyadic_square}
    \lean{TNLean.PEPS.AreaLaw.Geometry.latticeDyadicCell}
    \leanok
    For $k\in\N$ and $z\in\Z^2$, let $C_{k,z}$ be the integer sites of
    $\lbrack 2^kz_1,2^k(z_1+1))\times\lbrack 2^kz_2,2^k(z_2+1))$.
\end{definition}

\begin{theorem}[Integer dyadic nesting]
    \label{thm:area_law_lattice_dyadic_nesting}
    \lean{TNLean.PEPS.AreaLaw.Geometry.dyadicAncestor_eq_cellIndex,
        TNLean.PEPS.AreaLaw.Geometry.mem_latticeDyadicCell,
        TNLean.PEPS.AreaLaw.Geometry.mem_latticeDyadicCell_iff_mem_cell,
        TNLean.PEPS.AreaLaw.Geometry.card_latticeDyadicCell,
        TNLean.PEPS.AreaLaw.Geometry.latticeDyadicCell_zero,
        TNLean.PEPS.AreaLaw.Geometry.latticeDyadicCell_subset_of_common_mem,
        TNLean.PEPS.AreaLaw.Geometry.latticeDyadicCell_subset_parent}
    \leanok
    \uses{def:area_law_lattice_dyadic_square,def:area_law_dyadic_cells}
    Membership is given by coordinatewise integer division:
    $x\in C_{k,z}$ if and only if $\lfloor x/2^k\rfloor=z$.
    Equivalently, $C_{k,z}=\Z^2\cap 2^k\bigl(z+[0,1)^2\bigr)$ consists of
    the integer points of the half-open cell of
    Definition~\ref{def:area_law_dyadic_cells} with origin $o=0$; this is
    why the same symbol is used for both.
    Thus $|C_{k,z}|=4^k$ and $C_{0,z}=\{z\}$.
    Two intersecting squares are nested in scale order. In particular,
    $C_{k,z}\subseteq C_{k+1,\lfloor z/2\rfloor}$.
\end{theorem}
\begin{proof}\leanok
    \uses{def:area_law_lattice_dyadic_square,thm:area_law_dyadic_refinement,
        thm:area_law_dyadic_indices}
    The floor description of membership in the real cell at origin zero gives
    the identification with its integer points. Refine a single index to unit scale. Integer ancestry agrees with the
    floor index at origin zero. Composing ancestry maps shows that all sites
    of a fine square have the same coarse index, which proves nesting.
\end{proof}

\begin{definition}[Capped maximal contained squares]
    \label{def:area_law_capped_dyadic_partition}
    \lean{TNLean.PEPS.AreaLaw.Geometry.cappedDyadicPartition}
    \leanok
    \uses{def:area_law_lattice_dyadic_square}
    For a finite set $S\subseteq\Z^2$ and $K\in\N$, let $\mathcal P_K(S)$
    consist of the pairs $(k,z)$ with $k\le K$, $C_{k,z}\subseteq S$, and
    either $k=K$ or $C_{k+1,\lfloor z/2\rfloor}\not\subseteq S$.
    This is the selection in the covering argument of
    \cite[Lemma~9.4]{OpenAI2026AreaLaw}.
\end{definition}

\begin{theorem}[Exact capped partition]
    \label{thm:area_law_capped_dyadic_partition}
    \lean{TNLean.PEPS.AreaLaw.Geometry.mem_cappedDyadicPartition,
        TNLean.PEPS.AreaLaw.Geometry.exists_mem_cappedDyadicPartition,
        TNLean.PEPS.AreaLaw.Geometry.biUnion_cappedDyadicPartition,
        TNLean.PEPS.AreaLaw.Geometry.pairwiseDisjoint_cappedDyadicPartition,
        TNLean.PEPS.AreaLaw.Geometry.sum_card_cappedDyadicPartition,
        TNLean.PEPS.AreaLaw.Geometry.sum_pow_cappedDyadicPartition,
        TNLean.PEPS.AreaLaw.Geometry.card_cappedDyadicPartition_at_cap_le,
        TNLean.PEPS.AreaLaw.Geometry.cappedDyadicPartition_empty,
        TNLean.PEPS.AreaLaw.Geometry.cappedDyadicPartition_zero}
    \leanok
    \uses{def:area_law_capped_dyadic_partition}
    The squares selected by $\mathcal P_K(S)$ partition $S$. In particular,
    \begin{align}
        \sum_{(k,z)\in\mathcal P_K(S)}|C_{k,z}|
                                                       & =\sum_{(k,z)\in\mathcal P_K(S)}4^k=|S|, \\
        4^K\bigl|\{(k,z)\in\mathcal P_K(S):k=K\}\bigr| & \le |S|.
    \end{align}
    The empty set has an empty partition, and
    $\mathcal P_0(S)=\{(0,x):x\in S\}$.
\end{theorem}
\begin{proof}\leanok
    \uses{thm:area_law_lattice_dyadic_nesting,def:area_law_capped_dyadic_partition}
    For each $x\in S$, the unit square $C_{0,x}=\{x\}$ is contained in $S$.
    Choose
    \begin{align}
        k & =\max\bigl\{j\le K:C_{j,\lfloor x/2^j\rfloor}\subseteq S\bigr\}.
    \end{align}
    If $k<K$, maximality shows that the parent
    $C_{k+1,\lfloor x/2^{k+1}\rfloor}$ is not contained in $S$, so
    $(k,\lfloor x/2^k\rfloor)\in\mathcal P_K(S)$. Conversely, every
    selected square is contained in $S$. Let $(k,u),(l,v)\in\mathcal P_K(S)$
    with $k<l$ and $C_{k,u}\cap C_{l,v}\neq\varnothing$. Nesting gives
    \begin{align}
        C_{k+1,\lfloor u/2\rfloor}\subseteq C_{l,v}\subseteq S,
    \end{align}
    and $k<l\le K$, contradicting the selection of $(k,u)$. At equal scales the common site forces equal indices.
    Counting the disjoint union gives the sum identity; restricting its
    nonnegative summands to scale $K$ gives the cap bound.
\end{proof}

\begin{definition}[Mixed dyadic squares]
    \label{def:area_law_mixed_dyadic_square}
    \lean{TNLean.PEPS.AreaLaw.Geometry.mixedDyadicIndices}
    \leanok
    \uses{def:area_law_lattice_dyadic_square}
    Let $M_k(S)$ be the indices $z$ for which $C_{k,z}$ meets both $S$
    and its complement.
\end{definition}

\begin{lemma}[Counting integer quotient intervals]
    \label{lem:area_law_quotient_interval_count}
    \lean{TNLean.PEPS.AreaLaw.Geometry.quotient_interval_card}
    \leanok
    Let $a,b,u\in\mathbb Z$ with $a\le b$ and $u>0$. Then
    \begin{align}
        u\,\#\bigl([\lfloor a/u\rfloor,\lfloor b/u\rfloor]\cap\mathbb Z\bigr)
        \le b-a+2u.
    \end{align}
\end{lemma}
\begin{proof}\leanok
    The interval contains $\lfloor b/u\rfloor-\lfloor a/u\rfloor+1$ integers.
    Apply $u\lfloor b/u\rfloor\le b$ and
    $a<u(\lfloor a/u\rfloor+1)$.
\end{proof}

\begin{theorem}[Mixed dyadic squares of an integer rectangle]
    \label{thm:area_law_rectangle_mixed_count}
    \lean{TNLean.PEPS.AreaLaw.IntRect.card_mixedDyadicIndices_toFinset_le}
    \leanok
    \uses{def:area_law_safe_rectangle,def:area_law_lattice_dyadic_square,
        def:area_law_mixed_dyadic_square,lem:area_law_quotient_interval_count}
    Let $Q=[x_0,x_1]\times[y_0,y_1]\cap\mathbb Z^2$ be a nonempty integer
    rectangle, let $s$ be its size, and put $u=2^k$ for $k\in\mathbb N$.
    If $M_k(Q)$ is the set of indices of dyadic squares meeting both $Q$
    and its complement, then
    \begin{align}
        u\,\#M_k(Q)\le4s+8u.
    \end{align}
    This is the boundary-line count used in
    \cite[proof of Proposition~9.5, lines~717--729]{OpenAI2026AreaLaw}.
\end{theorem}
\begin{proof}\leanok
    \uses{def:area_law_lattice_dyadic_square,def:area_law_mixed_dyadic_square,
        lem:area_law_quotient_interval_count}
    Set $X=[\lfloor x_0/u\rfloor,\lfloor x_1/u\rfloor]\cap\mathbb Z$ and
    $Y=[\lfloor y_0/u\rfloor,\lfloor y_1/u\rfloor]\cap\mathbb Z$.
    An inside point places a mixed square's index in $X\times Y$.
    An outside point in the same square crosses one of the four sides of $Q$.
    Monotonicity of integer division then forces the corresponding index
    coordinate to equal an endpoint of $X$ or $Y$. Consequently
    \begin{align}
        \#M_k(Q)\le2\#X+2\#Y.
    \end{align}
    The quotient-interval estimate gives
    $u\,\#M_k(Q)\le2(x_1-x_0)+2(y_1-y_0)+8u\le4s+8u$.
\end{proof}

\begin{theorem}[Mixed parents below the cap]
    \label{thm:area_law_capped_mixed_parents}
    \lean{TNLean.PEPS.AreaLaw.Geometry.mem_mixedDyadicIndices,
        TNLean.PEPS.AreaLaw.Geometry.parent_mem_mixedDyadicIndices,
        TNLean.PEPS.AreaLaw.Geometry.cappedDyadicPartition_subset_mixed_refinement,
        TNLean.PEPS.AreaLaw.Geometry.card_cappedDyadicPartition_below_cap_le}
    \leanok
    \uses{def:area_law_mixed_dyadic_square,def:area_law_capped_dyadic_partition}
    If $(k,z)\in\mathcal P_K(S)$ and $k<K$, then
    $\lfloor z/2\rfloor\in M_{k+1}(S)$. Consequently,
    \begin{align}
        \bigl|\{(l,z)\in\mathcal P_K(S):l=k\}\bigr|\le4|M_{k+1}(S)|.
    \end{align}
\end{theorem}
\begin{proof}\leanok
    \uses{thm:area_law_lattice_dyadic_nesting,def:area_law_capped_dyadic_partition,
        thm:area_law_dyadic_refinement}
    A selected square is nonempty and contained in $S$, so its parent has
    an inside site. Below the cap, selection supplies an outside site in
    the parent. All selected indices at scale $k$ therefore belong to the
    one-step refinement of $M_{k+1}(S)$, whose cardinality is $4|M_{k+1}(S)|$.
\end{proof}
